\scanpage{312}
\mspage{11}
la seule connaissance de $C_\bullet\in\mathrm{Ab}_\bullet$, au lieu
et place de $X_*$; il faut, pour \deleted{\unsure{la connaissance}}
décrire en \uline{termes de} $C_\bullet$ \inserted{certaines des} structures
qu'on peut \inserted{ainsi} avoir sur $F(X_*)$, en termes
de celle de $F$, de tenir compte de
structures supplémentaires convenables
de $C_\bullet$, qui pourront s'expliciter en se
rappelant que $C_\bullet$ provient d'un $X_*$. \deleted{\unsure{Les}}
L'étude systématique de ces structures \deleted{\unsure{fera}}
\deleted{\unsure{est}} \inserted{peut être considérée comme}
le contenu essentiel de nos réflexions.
Mais bien sûr, \deleted{\unsure{si les structures}}\inserted{\unsure{dans}} \deleted{\unsure{d'ison}}
\unsure{si on a} sur $F$ une structure de
$h$-module, i.e. $F:\mathrm{Ss}^o\longrightarrow\mathrm{Ab}_h$
(où $\mathrm{Ab}_h$ est la catégorie des $h$-modules),
\deleted{il suffit}\inserted{il est inutile} de se donner de structure
suppl. sur $C_\bullet$, car \deleted{\unsure{car}} \unsure{qu'en fait}
\deleted{de complexe de $h$-chaînes sur $C_\bullet$
(i.e. se donner $h\longrightarrow\operatorname{End}_{\mathrm{Ab}_\bullet}(C_\bullet)$) pour
récupérer la structure complète de}
\unsure{que|car} le $K_\bullet$ qui décrit $F$ est
maintenant dans $(\mathrm{Ab}_h)_\bullet$ i.e. que
\[
h\longrightarrow\operatorname{End}_{\mathrm{Ab}_\bullet}(K_\bullet)
\]
$h$ opère aussi sur $\operatorname{Hom}_{\mathrm{Ab}_\bullet}(C_\bullet,K_\bullet)$,
dont on récupère ainsi la structure
de $h$-module. En fait, il suffit,
au lieu de $C_\bullet$, de connaître
seulement \inserted{$C_{h\bullet}=$} $C_\bullet\otimes_{\mathbf Z}\unsure{h|k}$ (not $\mathbb C_{h*}(X_*)$),
car on aura

\deleted{\unsure{la seule connais.}}
\[
\left\{\begin{aligned}
F(X_*)&\simeq\operatorname{Hom}_{(\mathrm{Ab}_h)_\bullet}(C_{h\bullet},K_\bullet)\\
C_{h\bullet}&=\mathbb C_{h*}(X_*)=\mathbb C_*(X_*)\otimes_{\mathbf Z}\unsure{h|k}
\end{aligned}\right.
\]
Plus généralement, revenant au cas d'une
$B$ cat. additive générale stable par $\varprojlim$
et $F:(\mathrm{Ss})^o\longrightarrow B$ \inserted{\unsure{fct}} et $K_\bullet\in B_\bullet$ se correspondant
par la prop. 2, on peut encore
écrire
\[
F(X_*)\simeq\operatorname{Hom}^{\bullet}_{\bullet\deleted{\unsure{B}}}(\mathbb C_\bullet(X_*),K_\bullet)
\]
en définissant le deuxième mb comme
\uline{objet de $B$}, en définissant d'abord le
bicomplexe de $B$
\[
\mathcal Hom^{\bullet}_{\bullet}(C_\bullet,K_\bullet)=
\bigl(\operatorname{Hom}_B(C_i,K_j)\bigr)_{i,j}
\]
(donc $\mathcal Hom^i_j(C_\bullet,K_\bullet)=\operatorname{Hom}_B(C_i,K_j)$)
--- où l'objet $\operatorname{Hom}(A,M)$, pour $A\in\mathrm{Ab}$
et $M\in B$, se définit par
\[
\operatorname{Hom}_B(N,\operatorname{Hom}(A,M))\simeq
\operatorname{Hom}_{\mathrm{Ab}}(A,\operatorname{Hom}_B(N,M))
\]
(le foncteur \deleted{\unsure{étant représenté}}\inserted{du $2^{\text{ème}}$ mb} étant bien
représentable en $N$, pour $A,M$ fixés, \uline{grâce}
à l'hyp. faite sur $B$ d'être stable par $\varprojlim$).
On définit alors
\[
\operatorname{Hom}_\bullet(C_\bullet,\deleted{\unsure{K_\bullet}},K_\bullet)
\overset{\mathrm{def}}=Z^0(\mathcal Hom^{\bullet}_{\bullet}(C_\bullet,K_\bullet))
\]
\margin{cycles de degré $0$ pour degré total \unsure{cohomolog.}}
Si enfin $B$ est une cat. $h$-linéaire
($h$ commutatif) i.e. munie d'un hom. d'
anneaux $h\longrightarrow\operatorname{End}(\mathrm{id}_B)$
\margin{313}
